% ED26-BM-RAT2021-11-2 C01--C05: exact Elqayam2021 psychology treatment.
\chapter{Deontic reasoning in psychology}
\label{sec:p19-rationality-11-2}

Deontic psychology investigates reasoning with and about norms. It includes
detecting a rule violation, interpreting a permission or promise, attributing
what an agent should do, generating a new normative conclusion, and evaluating
the standards used to assess reasoning. These are not interchangeable tasks.
Elqayam's 2021 chapter follows their historical development from the selection
task through social pragmatics, cognitive variability and is--ought inference.
Its central argument is that conformity to a specified inferential standard,
successful pursuit of goals, and participation in a normative social practice
interact but must not be conflated
\cite[pp.~651--655]{Elqayam2021DeonticPsychology}.

\begin{hbdefinition}{def:p19-r112-domain}{Norms, reasoning and moral orientation}
A deontic expression concerns obligation, permission, prohibition or a related
normative standing: must, may, must not, ought, should not, and need not are
examples. A \emph{norm} specifies what is required, permitted or prohibited
under stated circumstances; it is not merely a report of what usually happens.
\emph{Reasoning with norms} draws conclusions relative to given norms.
\emph{Reasoning about norms} can examine their applicability, interpretation,
source, conflict or proposed adoption. Empirically endorsing a conclusion,
having an obligation, and performing the action are distinct.

The chapter's note~1 separates \emph{deontic} from \emph{deontological}.
Deontic vocabulary is used by both consequence-oriented moral judgments and
judgments relying on preexisting duties. Deontological morality, in the
chapter's terminology, is the latter orientation, not all reasoning containing
the word \emph{ought}. An ethical classification does not by itself identify
a fast or slow psychological process
\cite[pp.~651,655, note~1]{Elqayam2021DeonticPsychology}.
\end{hbdefinition}

The opening examples range across traffic regulations, games, law, religion,
education and medical practice. They illustrate such different normative
contents as visiting a relative, requiring evidence for a recommendation,
prohibiting smoking in a specified train, and allowing a calculator in an
examination. These are examples of normative language, not current legal or
clinical guidance. The cited early game-rule protests concern a behavioral
measure of sensitivity to a practice, not a child's mastery of an entire
deontic calculus. Elqayam also distinguishes a first-order question---how
people reason from deontic premises---from a second-order question---which
standards researchers think reasoning \emph{should} satisfy. The latter
requires justification as well as observation
\cite[p.~651]{Elqayam2021DeonticPsychology}.

\section{The selection task: violation, falsification and information}
\label{sec:p19-r112-selection}

For this section, \(p\) and \(q\) name interpreted antecedent and consequent
conditions, and \(\neg\) denotes denial of a condition. The chapter's abstract
example displays A, N, 8 and 4, with each card having a letter on one side and
a number on the other. A rule links A to 8. Under the declared goal of testing
an exceptionless material rule over these cards, a violation has \(p\) together
with \(\neg q\). The A card can conceal a non-8 number; the 4 card can conceal
A. The N card cannot be an A case, and the 8 card cannot conceal a non-8
consequence. Thus A and 4 are the required checks for that specified complete
test. This case analysis fixes the benchmark; it does not identify the
mechanism that produced a participant's selection.

The drinking-rule example has beer, soft drink, age 22 and age 16 as its visible
cases. The task asks a police officer to detect breaches of a stipulated rule
requiring a beer drinker to be over nineteen. Beer and age sixteen are the
relevant checks: a beer drinker might be too young, and a young person might
be drinking beer. An underage beer drinker violates the regulation; that
occurrence need not make the regulation cease to apply. In contrast, a genuine
counterexample to an asserted universal descriptive claim makes that claim
false. This normative-rule versus descriptive-generalization distinction is
part of the task, not an arithmetic difference between the card arrays.

Elqayam reports approximately ten percent classical-answer selections for
typical abstract tasks and 74 percent for the cited Griggs--Cox drinking task,
and invokes the later selection-task review for robustness. Those are her
source-level summaries, not newly reconstructed samples or a pooled effect
estimate. The original experiments' recruitment, exclusions, precise task
versions and intervals are needed for a stronger population comparison.
The task's age rule is a stipulated experimental example, not current law.
The chapter's Wason (1960) historical citation is not used here to invent an
original four-card sample
\cite[p.~652]{Elqayam2021DeonticPsychology}.

\subsection{The source's explanatory alternatives}

Cheng and Holyoak's \emph{pragmatic reasoning schemas} are generalized but
context-sensitive structures organized around goals and relationships to
goals, including permission, obligation and causation. Calling a response
schema-driven still leaves its acquisition, representation and selection to
be explained. The chapter records criticism that a schema label alone may
redescribe rather than explain facilitation.

Cosmides and Tooby's evolutionary account instead proposes specialized
information-processing mechanisms shaped by recurrent adaptive problems.
Social exchange involves benefits and costs; cheater detection identifies
taking a benefit without satisfying its required cost. In that description,
the conditional roles are benefit and cost. Those roles must not silently
carry over to every differently worded permission. A successful social-contract
response alone does not discriminate an evolved specialization from a learned
goal-sensitive schema or another process producing the same choice.

Manktelow and Over emphasize social roles and subjective utility. In the
chapter's room-tidying promise, let \(p\) now mean that the son tidies and
\(q\) that he is allowed to play. The mother's valued outcome is tidiness;
her unfavorable case is \(\neg p\) with \(q\), play without tidying. The
son's unfavorable case is \(p\) with \(\neg q\), fulfilled effort without
the promised permission. Elqayam reports corresponding perspective-sensitive
selections. These are role-specific breach targets, not two inconsistent
answers to one unchanged task. A conditional permission does not automatically
assert a material regularity about what the son actually does
\cite[pp.~652--653]{Elqayam2021DeonticPsychology}.

Oaksford and Chater propose probabilistic information acquisition rather than
exhaustive material falsification: the value of examining a card depends on
which hypothesis is being tested, the distribution of possible outcomes and
the inquiry objective. The information-gain record in
\cref{sec:p19-r45-search} states these assumptions and distinguishes them
from costs and downstream decision utility. Neither the modal choice of
participants nor the word \emph{rational} proves that this objective is the
one specified by an experiment. Conversely, a classical benchmark does not
make every evidence-acquisition goal a complete-verification task.

Finally, the chapter discusses relations among deontic operators apart from
particular exchange benefits. It reports substantial competence alongside
errors described by Bucciarelli and Johnson-Laird's \emph{naive deontics}
account as illusory inferences. These are psychological meaning and
representation proposals, not a new sound-and-complete calculus established
by participant frequencies. For an actual monadic formal benchmark compare
\cref{def:p08-sdl}; for conditional obligation compare
\cref{def:p08-dyadic-obligation}. Neither formal presentation supplies the
participant's natural-language interpretation automatically.

\section{Speech acts, implicatures and practical commitments}
\label{sec:p19-r112-speech}

A promise can create an obligation by being made in a suitable social context;
it is not only a description of an already existing state. Elqayam's
collaborative-writing example concerns a deadline commitment. The chapter
connects this performative role to conversational pragmatics and distinguishes
\emph{inducements}, promises and threats, from \emph{advice}, tips and warnings.
Control over a consequence, intended effects, perceived value, authority,
interpretation and emotion can all affect a speech act's force. A threatened
consequence controlled by the speaker differs from a warning about a danger
the speaker does not control. Feeling anger at a broken promise, predicting
an agent's readiness to act, and declaring a formal consequence are different
outcomes of that social setting
\cite[p.~653]{Elqayam2021DeonticPsychology}.

The tidying promise can invite an additional understanding: without tidying,
permission will not be given. That \emph{inverse} implication is additional
to the stated forward relation. The \emph{converse} instead links the
consequence back to the antecedent. Elqayam reports more biconditional-like
responses to promises and threats than to indicative rules, but not an
equivalent strengthening for weaker tips and warnings. A social implicature
does not become a material entailment merely because it is often accepted.

The chapter's ``pattern X'' retains an important qualification. In a
truth-status task, the both-true and both-false antecedent/consequent cases
are judged true, the true-antecedent/false-consequent case false, and the
false-antecedent/true-consequent case irrelevant. That last response makes
this a \emph{defective} biconditional pattern, not the two-valued biconditional
truth table. Moreover, the greater frequency of inverse than converse
interpretations does not transfer directly to a preference for denial of the
antecedent in a conditional-inference task. The source explicitly notes the
absence of that expected preference. A sentence-interpretation task and an
inference task need their own response measures
\cite[p.~653]{Elqayam2021DeonticPsychology}.

The utility-conditionals treatment in \cref{sec:p19-rationality-6-4} records
speaker, hearer, actor, target and perceived outcome values. It distinguishes
an observer's attribution of a recommendation from the observer's own
endorsement. This is precisely why predicting the son's anger or the mother's
breach judgment does not establish which participant ought to prevail.

\section{Variability: developmental, cultural and individual}
\label{sec:p19-r112-variability}

Elqayam considers a possible common core without declaring identical reasoning
across every population. Her cross-cultural comparison names German, British,
Chinese and Tongan contexts. Shared conditional interpretations coexist with
variation in locally accepted conduct and emotional reactions. In the
formulation-choice example, explicit threats were selected more often by
German participants, whereas an indirect promise formulation was preferred
by Tongan participants, with British responses intermediate. These are
particular speech-act tasks, not representative measures of every member of
four cultures. Note~2 expands WEIRD as Western, Educated, Industrialized,
Rich and Democratic; that label neither describes an individual's whole
cultural history nor guarantees comparable recruitment
\cite[pp.~653--655, note~2]{Elqayam2021DeonticPsychology}.

The chapter also compares early game-rule protests, permission judgments,
promise-keeping and authority-sensitive judgments. Their reported emergence
ages range from two to seven depending on the study. This disagreement is not
resolved by selecting the earliest number: tasks differ in what they require
the child to infer, recognize or express. The adult individual-difference
literature likewise yields mixed relations between cognitive ability and
deontic selection performance. A task with strong contextual support need not
have the same score distribution or predictive relation as a demanding abstract
task. Neither early success nor a weak ability correlation establishes a
universal, innate and culturally invariant module
\cite[p.~654]{Elqayam2021DeonticPsychology}.

\subsection{An original reduced-array study with preschool children}
\label{sec:p19-r112-cummins}

Cummins's 1996 paper used rubber mice whose squeaking was revealed by
squeezing, distributed inside a model house and outside in a yard. The
descriptive rule said that squeaky mice were inside; the deontic rule required
them to remain inside. Visible location supplied a consequence-side case;
squeezing supplied the hidden antecedent information. Selecting outside mice
could discover a squeaky counterexample or violator. This \emph{reduced-array}
task presents a choice of location groups, not the full four-card array.
Children first pointed and verbally indicated which group to check, then
squeezed the mice. The measure was choosing the outside group, not explaining
deontic logic or demonstrating real-world compliance
\cite[pp.~824--826]{Cummins1996PreschoolDeontic}.

Experiment~1 recruited 44 children in the nominal three-year group and 40 in
the four-year group from Tucson day-care centers and preschools, with parental
consent. The younger deontic group ranged from 3 years 1 month to 3 years
11 months, median 3 years 6 months; its indicative counterpart ranged from
2 years 11 months to 3 years 11 months, median 3 years 5 months. Four-year
groups ranged from 4 years to 4 years 11 months, medians 4 years 4 months
and 4 years 3 months. Sessions were individual and videotaped, with house
location and query order counterbalanced. The deontic story introduced a cat
endangering noisy mice and a Queen Mouse authority; the indicative story asked
whether the experimenter was tricking the child. The full story contrast,
not the modal word alone, was the manipulation. Random assignment is not
stated in this method.

In the original-record subsections, \(p\) denotes a reported statistical
p-value, not the antecedent condition used in the selection discussion;
\(G^2\) is the paper's likelihood-ratio statistic. Outside-group choices
were 15/22 versus 7/22 in the three-year deontic/indicative groups,
and 17/20 versus 7/20 in the four-year groups. The reported percentages
were 68 versus 32, and 85 versus 35. The corresponding tests were
\(G^2(1)=5.95,p<.02\) and \(G^2(1)=11.03,p<.01\).
The paper supplies no confidence intervals for these proportions
\cite[pp.~824--825, table~1]{Cummins1996PreschoolDeontic}.

Experiment~2 recruited 70 nominal three-year-olds and 40 four-year-olds
from the same kinds of Tucson settings. The younger deontic group ranged
from 2 years 6 months to 3 years 11 months, median 3 years 6 months;
the indicative group from 2 years 11 months to 3 years 11 months, median
3 years 5 months. Both older groups ranged from 4 years to 4 years
11 months, median 4 years 6 months. The younger label therefore is not an
exact lower age bound. Counterbalancing was uneven among younger children
because of class sizes. Sessions were audiotaped.

The modified indicative story now included the cat, and a possibly wrong
Minnie Mouse assertion replaced the experimenter's deception framing.
One younger deontic participant pointed outside but squeezed all mice and
was excluded. Table~2 gives outside selections of 21/34 versus 13/35 for
the younger groups and 16/20 versus 6/20 for the older groups; reported
percentages were 62 versus 37 and 80 versus 30. Tests were
\(G^2(1)=4.23,p<.05\) and \(G^2(1)=10.60,p<.01\).
These are 109 analyzed children from 110 recruits, not 110 independently
successful tests, and not the same cohort as Experiment~1
\cite[p.~826, table~2]{Cummins1996PreschoolDeontic}.

The repetition with revised stories constrains the explanation that only
fear for the mice produced the difference. It does not isolate every
authority, linguistic or instruction effect, establish an infant capacity,
or measure general deductive competence. Cummins explicitly acknowledges
learning from frequent deontic experiences and analogy to familiar situations
as alternatives to innateness. Her evolutionary discussion favors a broader
deontic specialization related to dominance and reciprocal obligations rather
than exclusively human social exchange. Those are theoretical interpretations;
these two small, local, cross-sectional experiments cannot by themselves prove
the proposed evolutionary genealogy. Their internal conceptual replication
is not an independent cross-cultural replication
\cite[pp.~826--828]{Cummins1996PreschoolDeontic}.

\section{Generating norms and the is--ought problem}
\label{sec:p19-r112-is-ought}

The chapter's cooking example describes Ben wanting a better-tasting dish
to impress his girlfriend and an olive oil expected to improve its flavor.
It asks how reasoners arrive at the new conclusion that Ben should use
the oil. That conclusion adds normative standing to descriptive, goal-laden
information. The source calls the process \emph{deontic introduction}.
The recommendation can be withdrawn if the action--outcome link is undermined,
a goal conflicts, or another norm opposes the action, as in the example
of exploitative production. These are defeasible psychological transitions,
not proof that a desire by itself creates an overriding duty
\cite[p.~654]{Elqayam2021DeonticPsychology}.

\begin{hbexample}{def:p19-r112-bridge}{A successful recommendation is not a purely descriptive entailment}
Keeping the same cooking facts, one agent may regard improved flavor as
sufficient reason to use the oil; another may apply a conflicting purchasing
norm. Both can accept the descriptive action--flavor relation while giving
different normative judgments. A recommendation therefore requires its
goal, alternatives, costs, applicable norms and defeaters to be made explicit.
If the oil cannot improve the flavor, the stated instrumental reason fails;
if it can, that reason can still be overridden. These cases explain the
source's distinction without deciding a real purchasing or ethical question.
\end{hbexample}

Elqayam connects this issue with Hume's is--ought problem and Schurz's logical
treatment. The relevant boundary is a \emph{contentful normative conclusion
from the specified descriptive premises alone}, not an unqualified theorem
that no formula containing a deontic symbol can ever be logically valid.
A normative premise or bridge changes the argument; so does changing its
semantic standard. Note~3 identifies the naturalistic-fallacy issue concerning
natural properties as a related special case in the chapter's terminology.
This reference does not provide a full solution of that philosophical debate
\cite[pp.~654--655, note~3]{Elqayam2021DeonticPsychology}.

\subsection{An original moral-inference study and its qualifications}
\label{sec:p19-r112-moral-record}

Elqayam and colleagues' 2017 paper operationalizes \emph{exclusive coherence}
as a selective response to defeaters, not a mathematical coherence theorem
or a claim about every moral agent. Three vignettes were repeatedly used
in four tasks: ratings of utilitarian-action inference, deontological-action
inference, utilitarian moral rightness and deontological moral rightness.
Inference-strength ratings ran from one to seven; moral-rightness ratings
used a different one-to-seven question. Each task had one practice item and
three randomized test items, and the four task blocks were randomized
within participants. The same three vignettes occurring four times are
dependent observations, not twelve separately recruited people.

One vignette contrasts improving a religious group's diet with respecting
its beliefs. The causal premise predicts reduced vitamin deficiency; separate
questions concern the proposed dietary action and respect for the beliefs.
The method assesses six deontic expressions: must, must not, should, should
not, may and need not. The Experiment~1 design paragraph also lists going-to
and not-going-to despite specifying a six-level factor; the tables and
reported inference analyses use the six deontic operators. No eighth
operator condition is invented from that inconsistent enumeration
\cite[pp.~5--10, table~1]{ElqayamEtAl2017MoralInference}.

Experiment~1 recruited 116 paid CrowdFlower respondents selected for a good
quality-control history. Twenty-one were excluded for the reported eligibility,
consent or validation reasons, leaving 95 in four cells of sizes 25, 18,
28 and 24. The retained mean age reported in the age-comparison analysis
was 37.2 years versus 28.1 for excluded respondents; its test used only
\(69\) degrees of freedom, so the comparison does not supply a complete
age record for all recruits. Non-native English speakers and respondents
reporting dyslexia were among exclusions; this selected online population is
not representative of all adults. Normative conflict and negative outcome
information were randomly assigned between participants. The former invoked
a conflicting professional code; the latter added a severe adverse outcome
to the same action. These were fictional experimental materials, not medical
or professional advice.

The interaction of moral-judgment type with goal conflict was
\(F(1,91)=10.2,p=.002,\eta_p^2=.101\). For normative conflict the
corresponding interaction was \(F(1,91)=3.8,p=.054,\eta_p^2=.040\),
not below .05. Inference ratings showed operator/type/goal and
operator/type/normative-conflict interactions, with reported
\(F(2,137)=16.5,p<.001,\eta_p^2=.153\) and
\(F(2,137)=10.3,p<.001,\eta_p^2=.101\), respectively.
The paper used Greenhouse--Geisser adjustments where necessary and
Bonferroni comparisons. Effects were not uniform across operators:
permission expressed by \emph{may} was comparatively insensitive.
Figure error bars were one standard error, not 95 percent confidence
intervals; the paper's main interactions do not carry printed interval estimates
\cite[pp.~8--10, figs.~2--5]{ElqayamEtAl2017MoralInference}.

Experiment~2 recruited 235 paid CrowdFlower respondents, excluded 47 using
the reported criteria, and retained 188 in six cells of sizes 29, 34, 33,
25, 34 and 33. It retained the normative-conflict factor but replaced negative
outcome information with causal enhancement, disablers or alternative causes.
The diet vignette's enhancement makes the action effective and necessary;
its disabler prevents the stated outcome despite the action; its alternative
provides another route to the valued outcome without violating the stated
belief. A separate pretest of 71 respondents rated necessity and sufficiency.
Its complete materials and means are referred to supplementary Appendix~B;
the main article alone is not the complete pretest record.

The reported moral-judgment/type--causal interaction was
\(F(2,154)=3.78,p=.025,\eta_p^2=.047\); its type--normative-conflict
interaction was \(F(1,154)=10.8,p=.001,\eta_p^2=.065\).
Figure~7 prints utilitarian rightness means of 4.5 and 3.7 in low and high
normative conflict, and corresponding deontological means of 4.9 and 5.2.
Thus deontological judgments also changed under this manipulation, in the
opposite direction. The operator/type/causal interaction in inference strength
was \(F(4,278)=5.1,p=.001,\eta_p^2=.062\). Disabler comparisons
included an effect on deontological \emph{should}; alternative-cause effects
were less robust under the comparison correction. The source's overall
selectivity must retain these exceptions
\cite[pp.~10--14, figs.~6--9]{ElqayamEtAl2017MoralInference}.

The Experiment~2 inferential degrees of freedom and follow-up tests do not
give the 188-person methods denominator in every analysis, and the main
article does not supply an outcome-specific missing-data ledger. Its note~4
prints a positive contrast estimate 1.22 with a wholly negative 95 percent
interval \([-1.71,-.72]\); this sign inconsistency cannot be silently used
as a verified directional estimate. The original p.12 image confirms the
printed conflict. Neither result is reanalyzed here. Supplementary Appendices
A and B, needed for the complete item and pretest record, are not supplied
by the held main-paper PDF; their unavailable details are not fabricated.

The authors explicitly describe the selectivity as imperfect and decline to
infer ethical superiority from the findings. The two experiments conceptually
repeat the response pattern with different defeaters in a restricted item
family and online sampling frame; they are not independent multi-population
replications. Separate ratings reduce a forced-choice action/inaction confound,
but do not prove statistically independent psychological systems or a uniquely
identified causal processing chain. Nonsignificance in a deontological
comparison is not evidence of exactly zero effect or population equivalence.
Actual action and long-term norm adherence were not measured
\cite[pp.~14--16]{ElqayamEtAl2017MoralInference}.

The paper also describes the two courses of action as opposed. Even if they
are incompatible in a vignette, permission for one does not automatically
make the other forbidden: a rule can permit choosing either of two mutually
exclusive options. The contrary inference would need an additional normative
restriction, not only physical incompatibility. The handbook therefore does
not promote that sentence in the article's discussion into a general deontic
law; compare the explicit permission/obligation standing of
\cref{def:p08-sdl}.

\section{Normative benchmarks and the source's rationality debate}
\label{sec:p19-r112-debate}

Elqayam and Evans's descriptivist objection, reconstructed by the chapter,
concerns selecting a normative theory from behavioral evidence. One defense
of classical selection uses its association with cognitive ability and
motivation; a probabilistic defense can invoke modal selections. Either
move needs a premise connecting that observation to the proposed standard.
Being popular, or associated with a selected ability measure, does not
alone entail what someone ought to do. The chapter argues that using such
observations to settle a disputed norm risks reproducing an is--ought inference.
It records replies from Oaksford, Chater and Stanovich rather than presenting
that criticism as an uncontested impossibility theorem
\cite[pp.~654--655]{Elqayam2021DeonticPsychology}.

This does not prevent a researcher from explicitly choosing a benchmark and
asking whether a response meets it. The choice and the empirical question
must remain distinct. Complete violation checking follows from the declared
universal-rule task; efficient information acquisition requires a hypothesis
space and distribution; enforcing a conditional permission requires the
norm and the relevant authority/role; recommending an action requires its
goals, options, consequences and normative bridge. A norm does not become
binding merely because its representation predicts many responses.

\begin{hbmethod}{def:p19-r112-lookup}{Select a deontic task before judging a response}
Identify who speaks, who is addressed, who acts, whose outcome is valued,
and what authority or practice gives the rule its standing. Specify whether
the task tests descriptive truth, detects a violation of an accepted norm,
interprets a speech act, predicts behavior, attributes an attitude, endorses
a recommendation or judges morality. Declare the benchmark and premise
interpretation, including conditional scope and possible defeaters. Preserve
the actual response format: card selections, truth/irrelevance classifications,
inference ratings, rightness ratings and actions are different measurements.
Record recruitment, ages, language, exclusions, order, participant/item
dependence, scoring and uncertainty under
\cref{sec:p19-evidence-methods}.

A success case is detecting the young beer drinker under the stipulated
regulation, or distinguishing the mother's and son's different breach targets.
Failure cases include treating a detected violation as falsification of the
norm, equating permission with actual conduct, treating a pragmatic converse
as unchanged material entailment, or turning frequent recommendation ratings
into ethical validity. Separate generation of a new norm from inference
under a norm already assumed.
\end{hbmethod}

The chapter closes with specific research questions: broader cultural and
developmental comparisons, relations to moral judgment, whether endorsement
predicts action or reported readiness, and the neglected middle-strength
expressions \emph{ought} and \emph{should not}. Its final example is scientific
rigor itself: standards of evidence and research conduct have normative as
well as descriptive content. These are gaps and proposed research directions,
not an assertion that no relevant action studies exist. The original 2017
study supplies one qualified moral-inference interface, not the missing
general endorsement-to-action relation
\cite[p.~655]{Elqayam2021DeonticPsychology}.

\begin{proofboundary}{def:p19-r112-boundary}{Evidence and formal scope}
The complete 2021 chapter supplies the attributed historical, pragmatic and
methodological argument, including all three notes. The original preschool
paper and moral paper supply bounded task evidence; the moral paper's
supplementary item/pretest record and the other selected original-study
methods remain distinct dependencies. Review-level summaries do not acquire
invented sample sizes or intervals. No new mathematical axiomatization of
human cognition, complete empirical universality claim, legal advice,
clinical instruction or ethical ranking is established. The formal deontic
systems and their actual derivations remain separate from behavioral
endorsement and from the perceived social force of an utterance.
\end{proofboundary}
